\documentclass[10pt,a4paper]{amsart}
\usepackage[margin=19mm]{geometry}
\usepackage{amsmath,amssymb,mathtools}
\usepackage[scr=rsfs,cal=euler]{mathalfa}
\usepackage{xfrac}
\usepackage[unicode,hidelinks,bookmarks=false]{hyperref}
\newcommand{\ie}{i.e.\ }
\newcommand{\cf}{cf.\ }
\newcommand{\resp}{respectively\ }
\newcommand{\Subsection}{\subsection}
\providecommand{\nobreakdash}{\nobreak\hbox{-}}
\setlength{\emergencystretch}{2em}
\multlinegap0pt
\usepackage{amssymb}
\usepackage{xcolor}
\usepackage{tikz}
\newtheorem{prop}{Proposition}
\title{Hyperspherical Trigonometry and Corresponding Elliptic Functions}
\author{Paul Jennings, Frank Nijhoff}
\date{Source: arXiv:2211.13983v2 (CC BY 4.0)}
\begin{document}
\maketitle

\noindent\emph{Source and licence.} Hyperspherical Trigonometry and Corresponding Elliptic Functions. Version arXiv:2211.13983v2 (CC BY 4.0), distributed by arXiv under the Creative Commons Attribution 4.0 International licence, http://creativecommons.org/licenses/by/4.0/, as recorded in the arXiv licence metadata for this exact version and shown on the abstract page https://arxiv.org/abs/2211.13983v2. Full licence notice: \texttt{licenses/arxiv-2211.13983-CC-BY-4.0.txt}.\par\smallskip

\noindent\textit{Editorial scope.} Counted unit: the article's prop prop (source0.tex lines 1020--1028). The article's complete original proof of it is delivered with it (source0.tex lines 1030--1071, one \texttt{proof} environment of 42 lines). No mathematics is written, repaired or re-derived: the statement and the proof are the article's own lines, in the article's own notation, and no retained line is substituted or re-flowed.  No further source-local context is needed: every cross-reference of every retained line resolves inside the two delivered spans.  Nothing was omitted from any retained span, so the package carries no omission record.  No retained line contains a citation, so the wrapper carries no bibliography and no bibliography entry.  No retained line contains a figure environment, an image-inclusion command or drawing code, so no image is delivered and the package declares no asset dependency.  Editorial formatting only, in addition: an AMS \texttt{amsart} preamble that restates the article's own declarations and macros used by the retained lines, namely the environments \texttt{prop} on separate counters, together with the standard packages those lines need.  The retained environments print in this document as Proposition 1 in order of appearance, whereas the article prints the same items with the numbering of its own full document.  The complete native source of the article as distributed in the arXiv e-print for this exact version is bundled unchanged as \texttt{sources/133987/source0.tex}.
\begin{prop}
\begin{equation}
\begin{split}
\frac{\cos\phi_{ij}\cos\phi_{kl}-\cos\phi_{ik}\cos\phi_{jl}}{\cos\theta_{ij}\cos\theta_{kl}-\cos\theta_{ik}\cos\theta_{jl}}&=\frac{\cos\phi_{ik}\cos\phi_{jl}-\cos\phi_{il}\cos\phi_{jk}}{\cos\theta_{ik}\cos\theta_{jl}-\cos\theta_{il}\cos\theta_{jk}},\\
&=\frac{\cos\phi_{il}\cos\phi_{jk}-\cos\phi_{ij}\cos\phi_{kl}}{\cos\theta_{il}\cos\theta_{jk}-\cos\theta_{ij}\cos\theta_{kl}},\\
&=k_H.
\end{split}
\end{equation}
\end{prop}

\begin{proof}
Consider the numerator
\begin{equation}
\cos\phi_{ij}\cos\phi_{kl}-\cos\phi_{ik}\cos\phi_{jl},
\end{equation}
and rewrite this in terms of the central angles using the cosine rule,
\begin{equation}
\begin{split}
\cos\phi_{ij}\cos\phi_{kl}&-\cos\phi_{ik}\cos\phi_{jl}=
\frac{\left|\begin{array}{ccc}
\cos\theta_{ik} & \cos\theta_{jk} & \cos\theta_{kl}\\
1 & \cos\theta_{ij} & \cos\theta_{il}\\
\cos\theta_{ij} & 1 & \cos\theta_{jl}\\
\end{array}\right|}
{\sin(\theta_{ij},\theta_{ik},\theta_{jk})\sin(\theta_{ij},\theta_{il},\theta_{jl})}
\frac{\left|\begin{array}{ccc}
\cos\theta_{ik} & \cos\theta_{il} & \cos\theta_{ij}\\
1 & \cos\theta_{kl} & \cos\theta_{jk}\\
\cos\theta_{kl} & 1 & \cos\theta_{jl}\\
\end{array}\right|}
{\sin(\theta_{ik},\theta_{il},\theta_{kl})\sin(\theta_{jk},\theta_{jl},\theta_{kl})}\\
&-
\frac{\left|\begin{array}{ccc}
\cos\theta_{ij} & \cos\theta_{jk} & \cos\theta_{jl}\\
1 & \cos\theta_{ik} & \cos\theta_{il}\\
\cos\theta_{ik} & 1 & \cos\theta_{kl}\\
\end{array}\right|}
{\sin(\theta_{ij},\theta_{ik},\theta_{jk})\sin(\theta_{ik},\theta_{il},\theta_{jk})}
\frac{\left|\begin{array}{ccc}
\cos\theta_{ij} & \cos\theta_{il} & \cos\theta_{ik}\\
1 & \cos\theta_{jl} & \cos\theta_{jk}\\
\cos\theta_{jl} & 1 & \cos\theta_{kl}\\
\end{array}\right|}
{\sin(\theta_{ij},\theta_{il},\theta_{jl})\sin(\theta_{jk},\theta_{jl},\theta_{kl})}.
\end{split}
\end{equation}
Multiplying this out and factorising, this reduces to
\begin{equation}
\cos\phi_{ij}\cos\phi_{kl}-\cos\phi_{ik}\cos\phi_{jl}=(\cos\theta_{ij}\cos\theta_{kl}-\cos\theta_{ik}\cos\theta_{jl})k_H,
\end{equation}
from which the result follows.
\end{proof}

\end{document}
